\subsection{Policy content coding}
\label{sec:si-policy-methods}

The policy evidence is represented at three linked levels: clause, document
and city.  The research-window corpus contains 601 locally applicable
documents covering 176 cities.  Of these, 207 documents covering 88 cities
contain 466 retained requirement clauses.  A document enters the analytical
content corpus when it has a local administrative scope, an applicable date
within the study frame, and a requirement or condition that can be linked to
charging infrastructure or its operation.  A clause is the classification
unit; document and city identifiers are retained separately and are not
additive counts.  Original language, translated wording, document identity,
source locator, publication date, effective date and date precision are
retained for every evidence record.

The codebook records four functional groups: information and supervision;
operation and maintenance; continuing operation or payment; and assessable or
subsidy-linked conditions.  A label describes what a text requires, monitors
or conditions.  It is not a measure of compliance, policy intensity or realized
service.  Labels are non-exclusive: a clause may require both reporting and
maintenance, or condition payment on continued operation.  Functional city
counts therefore overlap and cannot be summed to a policy total.  Evidence is
located by the retained source fragment, document identifier and
clause identifier.  Because labels are non-exclusive, the functional counts
describe the number of clauses meeting each requirement definition rather than
a partition of documents or cities.

\subsection{Governance timing registry}
The timing registry is separate from the 88-city content scope and contains 95
cities with a qualifying governance record.  Construction support is coded
$B$, information or supervision $R$, and operating-service responsibility
$S$; $G$ denotes the combined path following the first qualifying $R$ or $S$
requirement.  Here $B$ identifies construction-support events; the annual
regressor $B_{i,t-1}$ defined below is a prior-year document count.
Publication and effective dates are retained separately.  For
each record, let $[a,b]$ be the inclusive interval implied by the earliest known
lower bound and latest known upper bound at the reported precision.  A calendar
quarter is clean pre-policy only when its end precedes $a$, transitional when it
intersects $[a,b]$, and fully post-policy only when its start follows $b$.
Year-only dates remain intervals and are not assigned an artificial event
quarter.  Unknown coverage or unresolved timing remains unknown and is not
coded as untreated.  A first year with observed facilities is not interpreted
as a construction or commissioning date.

The timing table retains the city-level dates, precision, pre-period support and
first/second-year availability needed by the governance comparison.  Content
coverage, governance timing, charging observations and their intersections
remain separate denominators in the accompanying overlap tables.  This keeps a
documented responsibility distinct from the later quarterly support and from
the annual association estimand.

\subsection{Annual construction-support year assignment}
For the annual policy count $B$, a document is assigned to a calendar year only
when its effective year is known.  When both effective and publication years are
known, the assigned year is the later of the two; when the effective year is
known but the publication year is not, the effective year is used.  When the
effective year is unknown, the assigned year remains unknown even if a
publication year is available.  The annual regressor $B_{i,t-1}$ is the number
of qualifying construction-support policy documents assigned to the preceding
calendar year $t-1$; it is a prior-year document count, not a cumulative stock.
The $R$ and $S$ indicators are already defined at the end of the preceding year
and are not lagged a second time.  This year-assignment rule is distinct from
the quarterly date-interval rule above: a quarter is cleanly pre-policy only
when its end precedes the lower bound, transitional when it intersects the
interval, and fully post-policy only when its start follows the upper bound.
